% ============================================================
% f2012_hubei_science_q12_flowchart.tex
% 2012 湖北卷（理科）第 12 题 程序框图
% 初值 a=1、s=0、n=1；循环体 s=s+a、a=a+2；判断 n<3，是则 n=n+1 返回，
% 否则输出 s。运行结果 s=9。
% 来源 PDF：普通高考/2012/2012湖北理.pdf（SHA-256 b84a88bc…07cc）
% 原页：第 1 页（A3 横排）第 3 栏，印面 737，第 12 题题干下方
% ============================================================
\begin{tikzpicture}[>=Stealth, line width=0.7pt, font=\small]
  % 开始
  \draw[rounded corners=8pt] (-0.55,0.50) rectangle (0.55,0.10);
  \node at (0,0.30) {开始};
  \draw[->] (0,0.10) -- (0,-0.30);

  % 初值
  \draw (-1.55,-0.30) rectangle (1.55,-0.80);
  \node at (0,-0.55) {$a=1,\ s=0,\ n=1$};
  \draw[->] (0,-0.80) -- (0,-1.20);

  % s = s + a
  \draw (-0.85,-1.20) rectangle (0.85,-1.70);
  \node at (0,-1.45) {$s=s+a$};
  \draw[->] (0,-1.70) -- (0,-2.05);

  % a = a + 2
  \draw (-0.85,-2.05) rectangle (0.85,-2.55);
  \node at (0,-2.30) {$a=a+2$};
  \draw[->] (0,-2.55) -- (0,-2.85);

  % 判断 n < 3
  \draw (0,-2.85) -- (1.05,-3.55) -- (0,-4.25) -- (-1.05,-3.55) -- cycle;
  \node at (0,-3.55) {$n<3$?};

  % 是：n = n + 1，回边
  \draw[->] (0,-4.25) -- (0,-4.60);
  \node[right] at (0.08,-4.38) {是};
  \draw (-0.80,-4.60) rectangle (0.80,-5.10);
  \node at (0,-4.85) {$n=n+1$};
  \draw (-0.80,-4.85) -- (-2.15,-4.85) -- (-2.15,-1.45);
  \draw[->] (-2.15,-1.45) -- (-0.85,-1.45);

  % 否：输出 s
  \draw[->] (1.05,-3.55) -- (2.15,-3.55);
  \node[above] at (1.30,-3.50) {否};
  \draw (2.45,-3.30) -- (3.95,-3.30) -- (3.45,-3.80) -- (1.95,-3.80) -- cycle;
  \node at (2.95,-3.55) {输出 $s$};
  \draw[->] (2.95,-3.80) -- (2.95,-4.50);

  % 结束
  \draw[rounded corners=8pt] (2.35,-4.50) rectangle (3.55,-4.95);
  \node at (2.95,-4.72) {结束};
\end{tikzpicture}
